% Sign Inventory
\section{Sign Inventory Overview}

The Indus script comprises approximately 417 distinct signs (following Mahadevan's concordance), though estimates vary depending on treatment of variants and damaged exemplars.

\subsection{Classification Criteria}

Signs are classified based on multiple criteria:

\subsubsection{Morphological Classification}
\begin{itemize}
    \item Geometric shapes (circles, squares, triangles)
    \item Linear forms (vertical strokes, horizontal bars)
    \item Anthropomorphic figures
    \item Zoomorphic figures
    \item Botanical forms
    \item Object representations (tools, vessels, structures)
    \item Abstract symbols
\end{itemize}

\subsubsection{Functional Classification}
Based on distributional analysis:
\begin{itemize}
    \item Instantiation signs (introduce entities)
    \item Counting signs (numerical values)
    \item Relational signs (connect entities)
    \item Qualification signs (add attributes)
    \item Terminal signs (close sequences)
\end{itemize}

\subsection{Sign Frequency Distribution}

The sign frequency follows a characteristic pattern:
\begin{itemize}
    \item Approximately 10 signs account for 50\% of all occurrences
    \item Approximately 40 signs account for 80\% of all occurrences
    \item Approximately 200 signs are singletons (appearing only once)
\end{itemize}

This distribution is typical of administrative/operational systems rather than phonetic writing systems.

\subsection{Positional Constraints}

Signs exhibit strong positional preferences:
\begin{itemize}
    \item Certain signs occur exclusively or predominantly in initial position
    \item Other signs occur primarily in terminal position
    \item Some signs are restricted to medial positions
    \item A small set of signs shows positional flexibility
\end{itemize}

These constraints are fundamental to the operational interpretation.
